\section{Conclusiones}

El desarrollo de este proyecto permitió el desarrollo de habilidades con relación a una introducción a los comandos de Bases de Datos SQL y una profundización del desarrollo del patrón Modelo-Vista-Controlador esto a través de las especificaciones dadas por el propio proyecto del reproductor de música.

La elección de tener una arquitectura ``híbrida'' con dos lenguajes de programación (Rust y Flutter) permitió el desglose de responsabilidades y además, fue una decisión importante a la hora de equilibrar el rendimiento y el escribir la interfaz de usuario, esto ya que Rust, se caracteriza por la seguridad en memoria, y rapidez, esto por su sistema de ``ownership'', el cual se explicó en el proyecto pasado, y por su parte Flutter, al ser un toolkit destinado a la creación de interfaces gráficas multiplataforma, resultó poder crear la interfaz de una forma un poco más clara, además de poder ser reactiva y modular, esto para poder llevar acabo de buena forma el patron Modelo-Vista-Controlador, y tener un modelo basado en funcionalidades como se describió anteriormente, mientras que utilizar la herramienta de Flutter\_Rust\_Bridge, mostró ser una de las opciones más óptimas para poder integrar estos dos lenguajes de programación, esto gracias a que es posible invocar funciones de Rust de forma asíncrona sin comprometer la fluidez de la interfaz gráfica. 

Asimismo, la separación del sistema en tres módulos independientes (\texttt{minero}, \texttt{nucleo\_dao} y \texttt{reproductor}) aportó ventajas clave en mantenimiento y pruebas:
\begin{itemize}
    \item El encapsulamiento de las consultas en el patrón \textbf{DAO} evitó la duplicación de código SQL y permitió validar cada operación de persistencia mediante pruebas unitarias automatizadas (\texttt{cargo test}) antes de integrarla con la interfaz.
    \item El diseño del \textbf{minero} como componente autónomo con salida tabular en consola (\textit{visualizer}) posibilitó depurar la extracción de metadatos de manera aislada y cumplir con el requisito de contar con ejecutables independientes.
    \item La edición de registros directamente sobre la base de datos SQLite aseguró el cumplimiento de la restricción no destructiva del sistema, permitiendo corregir información de canciones, artistas, grupos y relaciones sin alterar los archivos de audio originales.
    \item El uso de rutas estándar dependientes del sistema operativo y la incrustación del esquema SQL dentro del binario garantizaron la portabilidad de los ejecutables tanto en \textbf{macOS} como en \textbf{Linux}.
\end{itemize}

Finalmente, como trabajo futuro y oportunidades de mejora se identifica la refactorización de las clases de dominio compartidas (como la clase \texttt{Rola}) hacia una capa superior en Flutter para reducir el acoplamiento entre controladores, la ampliación del minero para soportar formatos adicionales sin pérdida (como \texttt{FLAC} o \texttt{WAV}) y la incorporación de listas de reproducción personalizadas persistentes en la base de datos.